\documentclass[11pt]{amsart}
\usepackage[T1]{fontenc}
\usepackage{lmodern,amsmath,amssymb,amsthm,microtype}
\usepackage[margin=1.15in]{geometry}
\usepackage[colorlinks=true,linkcolor=blue,citecolor=blue,urlcolor=blue]{hyperref}
% Repository locator for the public (non-anonymous) build of
% non_mf_groups_exist.tex.  Kept in a separate file so that the anonymous
% build has an anonymous source archive and not merely an anonymous PDF:
% exclude this file from a double-blind submission and the manuscript still
% compiles, printing "supplied with this submission as supplementary
% material" in place of the link.
%
% Loaded only when \anonymousfalse is set in the preamble.
\newcommand{\coderepository}{https://github.com/SauersML/group-approximation}
\newcommand{\coderepositoryname}{github.com/SauersML/group-approximation}

\hypersetup{pdftitle={An infinite simple Kazhdan sofic group},pdfauthor={\paperauthorname}}
\newtheorem{theorem}{Theorem}
\newtheorem{corollary}[theorem]{Corollary}
\newcommand{\F}{\mathbb F}
\newcommand{\Z}{\mathbb Z}
\newcommand{\LC}{\operatorname{LC}}
\newcommand{\EL}{\operatorname{EL}}
\newcommand{\GL}{\operatorname{GL}}
\newcommand{\SL}{\operatorname{SL}}
\newcommand{\PSL}{\operatorname{PSL}}
\newcommand{\doi}[1]{\href{https://doi.org/#1}{\nolinkurl{doi:#1}}}
\raggedbottom
\emergencystretch=1.5em

\begin{document}
\title{An infinite simple Kazhdan sofic group}
\author{\paperauthorname}
\date{}
\subjclass[2020]{Primary 20E32; Secondary 22D10, 37B10, 16S35}

\begin{abstract}
For every infinite minimal subshift $X$, the group
$\EL_3(\LC(X,\F_2)\rtimes\Z)$ is infinite, simple, finitely generated,
and has property~\textup{(T)}. It is locally embeddable into finite
groups, hence sofic and hyperlinear. This answers the questions of
Brown, Ozawa, and Pestov. The proof combines periodic approximation
with a commutator supported on a finite tower. Every Turing degree
occurs as the word-problem degree of one of these groups.
\end{abstract}
\maketitle

Can an infinite simple Kazhdan group be hyperlinear? Brown asked this
in its von Neumann algebra form in 2001~\cite[\S11, Question~7]{Brown},
and Ozawa in the hyperlinear form in 2003~\cite[p.~527]{Ozawa}.
Pestov's Open question~9.1 adds the sofic version~\cite{Pestov}.
The following construction answers both forms positively.

\begin{theorem}\label{thm:main}
Let $X\subseteq A^{\Z}$ be an infinite minimal subshift over a finite
alphabet, with shift $T$. Then
\[
  G_X=\EL_3\bigl(\LC(X,\F_2)\rtimes_T\Z\bigr)
\]
is an infinite finitely generated simple group with Kazhdan's
property~\textup{(T)} that is locally embeddable into finite groups
\textup{(LEF)}. In particular, it is sofic and hyperlinear.
\end{theorem}

The subshift algebra is classical; see Nekrashevych~\cite{Nekrashevych}.
Property~(T) is supplied by Ershov--Jaikin-Zapirain~\cite{EJZ}.
Our finite models use the periodic-approximation mechanism of
Grigorchuk--Medynets~\cite{GM}. The new step is the simplicity argument:
a nontrivial normal subgroup meets a finite simple matrix group
supported on a clopen tower, and therefore contains an elementary
matrix. This is an explicit elementary-matrix extraction argument,
in the normal-structure tradition described by Stepanov~\cite{Stepanov}.
Thom had already constructed a finitely generated Kazhdan LEF group
that is not residually finite~\cite[Theorem~1.4]{Thom}; his example
is not simple.

\section{Proof of Theorem~\ref{thm:main}}

\subsection*{The ring and property (T)}
Write $e_U$ for the indicator of a clopen set $U\subseteq X$. Our
conventions are $(Tx)_n=x_{n+1}$ and
\[
 R=\LC(X,\F_2)\rtimes_T\Z
   =\left\{\sum_j f_j u^j:\text{finite sums}\right\},
 \qquad ufu^{-1}=f\circ T^{-1},
\]
where $\LC$ denotes locally constant functions. Thus $ue_Uu^{-1}=e_{TU}$.
Minimality and infiniteness imply that $T$ has no periodic points.
For each $L$, every point therefore has a clopen neighborhood $U$ with
$U\cap T^jU=\varnothing$ for $0<|j|\le L$. Compactness gives finite
clopen partitions with this property, refining any prescribed clopen
partition.

The ring $R$ is generated by $u,u^{-1}$ and the letter indicators
$e_a=e_{\{x:x_0=a\}}$: their translates and products give all cylinder
indicators. Put $G=\EL_3(R)$ and $e_{ij}(r)=I_3+rE_{ij}$, where $E_{ij}$
is the usual matrix unit. The elementary identities
\begin{equation}\label{eq:elementary}
 \begin{aligned}
 e_{ij}(r+s)&=e_{ij}(r)e_{ij}(s),\qquad i\ne j,\\
 [e_{ik}(r),e_{kj}(s)]&=e_{ij}(rs),\qquad i,j,k\text{ distinct}.
 \end{aligned}
\end{equation}
show that the matrices $e_{ij}(s)$, for
$s\in\{1,u,u^{-1}\}\cup\{e_a:a\in A\}$, generate $G$.
Here $[g,h]=ghg^{-1}h^{-1}$.
Ershov--Jaikin-Zapirain's theorem says that $\EL_n(R)$ has property~(T)
for every finitely generated unital associative ring $R$ and
$n\ge3$~\cite[Theorem~1.1]{EJZ}. Hence $G$ has property~(T).
It is infinite because $e_{12}(\LC(X,\F_2))$ is infinite.

We will also need that $R$ is simple and $Z(R)=\F_2$.
These are standard groupoid-algebra facts~\cite{BCFS,ClarkEdie,Steinberg};
here is the direct argument. A nonzero two-sided ideal contains an
element $r=\sum_j f_ju^j$ with $f_0\ne0$, after multiplication by a
power of $u$. Choose a nonempty clopen $U$ on which $f_0=1$, disjoint
from $T^jU$ for every nonzero exponent occurring in $r$. Then
$e_Ure_U=e_U$. Finitely many translates of $U$ cover $X$, so the ideal
contains $1=1-\prod_i(1-e_{T^{n_i}U})$.
If $r$ is central, commuting with all clopen indicators forces $f_j=0$
for $j\ne0$: otherwise a point in the support of $f_j$ and its
$T^{-j}$-translate can be separated by a clopen set. Commuting with
$u$ then makes $f_0$ invariant, hence constant by minimality.
Finally, commuting with every $e_{ij}(1)$ forces a central element of
$G$ to be $cI_3$; commuting with all $e_{ij}(r)$ gives
$c\in Z(R)^\times=\{1\}$. Thus $Z(G)=\{1\}$.

\subsection*{Finite models}
A group is LEF if every finite subset embeds injectively into a finite
group, preserving all products that stay in that subset.
We construct such models for $R$, and then apply them entrywise.

For each $k$, there is a periodic sequence $y=w^\infty$ with exactly
the same words of length $2k+1$ as $X$. Indeed, minimality implies
uniform recurrence. In a point of $X$, choose two occurrences of the
same word of length $2k$, separated by a segment $w$ long enough to
contain every word of length $2k+1$. Repeating this segment introduces
no new word of that length at the join. Its length $N$ can be chosen
arbitrarily large.

Let $P\delta_n=\delta_{n+1}$ on $\F_2^{\Z/N\Z}$. If $f$ depends on a
window $[-\rho,\rho]$, evaluate its local table along $y$ to obtain a
diagonal matrix $D_y(f)$. Send
\[
  \sum_j f_ju^j\ \longmapsto\ \sum_j D_y(f_j)P^j.
\]
For any fixed finite list of sums and products in $R$, choose $k$
larger than all windows involved, including those shifted in the
products. The relation
$P^iD_y(f)P^{-i}=D_y(f\circ T^{-i})$ then preserves every listed
operation. Nonzero coefficients remain nonzero, since every allowed
window occurs in $y$. Taking $N$ larger than twice the largest absolute
exponent makes the corresponding cyclic diagonals distinct; hence
nonzero differences also remain nonzero.

For a finite subset of $G$, include the entries of its elements,
their inverses, their differences, and all intermediate sums and
products needed for matrix multiplication. The resulting model is an
injective partial homomorphism into $\GL_{3N}(\F_2)$. Thus $G$ is LEF.
The regular permutation actions of these finite groups give sofic
models; their permutation matrices give hyperlinear models
\cite{Pestov}.

\subsection*{A finite tower detects every normal subgroup}
Suppose the levels $T^aU$ are disjoint for all indices needed below.
For $|a|,|b|\le m$ and clopen $W\subseteq U$, put
\[
  E_{ab}(W)=e_{T^aW}u^{a-b}.
\]
These satisfy
$E_{ab}(W)E_{cd}(W')=\delta_{bc}E_{ad}(W\cap W')$.
Their span is a subring
\[
  B_m(U)\cong M_{2m+1}(\LC(U,\F_2)),
\]
with unit $e_{\bigcup_{|a|\le m}T^aU}$. If $r,r'$ involve only powers
$u^j$ with $|j|\le w$, then
\begin{equation}\label{eq:absorb}
  r B_m(U)r'\subseteq B_{m+w}(U),
\end{equation}
provided the larger tower is disjoint. To see this, multiplying
$E_{ab}(W)$ on the left by $fu^i$ and on the right by $f'u^j$ changes
its indices to $(a+i,b-j)$ and restricts its clopen coefficient.

Let $1\ne N\trianglelefteq G$, and choose $1\ne g\in N$.
Let $w\ge0$ bound the absolute values of the exponents in all entries
of $g$ and $g^{-1}$.
Choose a finite clopen partition $\mathcal P$ refining the letter
partition, with $C\cap T^jC=\varnothing$ for $C\in\mathcal P$ and
$0<|j|\le2w+3$. Some matrix
\[
 h=e_{ij}(s),\qquad
 s\in\{e_C,e_Cu,e_Cu^{-1}:C\in\mathcal P\},
\]
does not commute with $g$. Otherwise, by~\eqref{eq:elementary}, the
coefficients $s$ for which $g$ commutes with every $e_{ij}(s)$ form
a subring containing all ring generators, so $g\in Z(G)=\{1\}$.

Each such $s$ lies in $B_1(U)$ for a translate $U$ of $C$:
\[
 e_C=E_{00}(C),\quad e_Cu=E_{10}(T^{-1}C),\quad
 e_Cu^{-1}=E_{-1,0}(TC).
\]
Set $k=[g,h]\in N\setminus\{1\}$ and $m=w+1$.
By~\eqref{eq:absorb}, both $k-I_3$ and $k^{-1}-I_3$ lie in
$M_3(B_m(U))$. If $e$ is the unit of $B_m(U)$, then
$k=(1-e)I_3+ek e$. Its corner part is invertible and therefore
corresponds to a locally constant function
\[
 \kappa:U\longrightarrow\GL_d(\F_2),\qquad d=3(2m+1)\ge9.
\]
Choose nonempty clopen $W\subseteq U$ on which $\kappa$ is a constant
$\kappa_0\ne I_d$.

Let $H_W$ be the copy of $\GL_d(\F_2)$ acting on the tower over $W$
and as the identity elsewhere. It lies in $G$. Indeed, index its
coordinates by $(p,a)$, where $1\le p\le3$ and $|a|\le m$.
A transvection between $(p,a)$ and $(q,b)$ with $p\ne q$ is
$e_{pq}(E_{ab}(W))$. For $p=q$ and $a\ne b$, take $q\ne p$ and use
\[
 [e_{pq}(E_{ab}(W)),e_{qp}(E_{bb}(W))]
       =I_3+E_{ab}(W)E_{pp}.
\]
This follows from $E_{ab}E_{bb}=E_{ab}$ and $E_{bb}E_{ab}=0$.
These transvections generate $H_W$.

Over $\F_2$, the group $\GL_d(\F_2)=\SL_d(\F_2)=\PSL_d(\F_2)$
is simple and centerless for $d\ge3$. Conjugation by $k$ on $H_W$
is conjugation by $\kappa_0$, so $[k,H_W]\subseteq N\cap H_W$
contains a nonidentity element. Since $N\cap H_W$ is normal in $H_W$,
we get $H_W\subseteq N$. In particular, $e_{12}(e_W)\in N$.

Finally, $I_N=\{r\in R:e_{12}(r)\in N\}$ is a two-sided ideal:
permutation matrices in $G$ move any elementary position to any other,
and~\eqref{eq:elementary} gives addition and multiplication on either
side by arbitrary ring elements. Since $0\ne e_W\in I_N$ and $R$ is
simple, $I_N=R$. Thus $N=G$, completing the proof.\hfill$\square$

\section{Every word-problem degree}

\begin{corollary}
Every Turing degree occurs as the word-problem degree of an infinite
simple Kazhdan LEF group. Consequently, there are continuum many
isomorphism classes of these groups.
\end{corollary}

\begin{proof}
The word problem of $G_X$, in the finite generators above, has the same
Turing degree as the language $L(X)$ of finite words occurring in $X$.
An oracle for $L(X)$ decides whether a matrix word equals $I_3$, by
checking the coefficient tables of its difference from $I_3$ on the
allowed windows. Conversely, for a word $v=v_0\cdots v_{n-1}$, the cylinder
indicator
\[
 e_{[v]}=\prod_{t=0}^{n-1}u^{-t}e_{v_t}u^t
\]
can be converted effectively, using~\eqref{eq:elementary}, to a group
word representing $e_{12}(e_{[v]})$. This word equals $1$ precisely
when $v\notin L(X)$.

Now use the infinite minimal Sturmian subshift of an irrational slope
$\alpha\in(0,1)$~\cite{MorseHedlund}. Its language is computable from
$\alpha$: words of length $n$ are determined by the cyclic order of
the distinct rotation endpoints $-j\alpha\bmod1$, $0\le j\le n$.
Conversely, the minimum number of $1$'s in a word of length $n$ is
$\lfloor n\alpha\rfloor$, so the language computes $\alpha$.
Every Turing degree has an irrational representative: use $\sqrt2-1$
for the computable degree, and interleave the characteristic sequence
of a noncomputable set with $1$'s in a binary expansion for any other
degree. There are continuum many degrees, and the word-problem degree
of a finitely generated group is an isomorphism invariant.
\end{proof}

\subsection*{Origin and authorship}
Claude (Anthropic) found the construction and original proofs under
the author's direction on September~12, 2026, and wrote the original
manuscript. Codex (OpenAI) shortened and edited this version.
The author is responsible for the final manuscript. Proof records are
in the \href{https://github.com/SauersML/group-approximation}{project repository}.

\begin{thebibliography}{99}

\bibitem{Brown}
N.~P. Brown, \emph{Tracial invariants, classification and
$\mathrm{II}_1$ factor representations of Popa algebras},
\href{https://arxiv.org/abs/math/0111286}{arXiv:math/0111286} (2001).

\bibitem{BCFS}
J.~H. Brown, L.~O. Clark, C.~Farthing, and A.~Sims,
\emph{Simplicity of algebras associated to \'etale groupoids},
Semigroup Forum \textbf{88} (2014), 433--452.
\href{https://arxiv.org/abs/1204.3127}{arXiv:1204.3127}.

\bibitem{ClarkEdie}
L.~O. Clark and C.~Edie-Michell,
\emph{Uniqueness theorems for Steinberg algebras},
Algebr. Represent. Theory \textbf{18} (2015), 907--916.
\href{https://arxiv.org/abs/1403.4684}{arXiv:1403.4684}.

\bibitem{EJZ}
M.~Ershov and A.~Jaikin-Zapirain,
\emph{Property~\textup{(T)} for noncommutative universal lattices},
Invent. Math. \textbf{179} (2010), 303--347.
\doi{10.1007/s00222-009-0218-2}.

\bibitem{GM}
R.~Grigorchuk and K.~Medynets,
\emph{On algebraic properties of topological full groups},
Sb. Math. \textbf{205} (2014), 843--861.
\doi{10.1070/SM2014v205n06ABEH004400}.

\bibitem{MorseHedlund}
M.~Morse and G.~A. Hedlund,
\emph{Symbolic dynamics II. Sturmian trajectories},
Amer. J. Math. \textbf{62} (1940), 1--42.
\doi{10.2307/2371431}.

\bibitem{Nekrashevych}
V.~Nekrashevych, \emph{Growth of \'etale groupoids and simple algebras},
Internat. J. Algebra Comput. \textbf{26} (2016), 375--397.
\href{https://arxiv.org/abs/1501.00722}{arXiv:1501.00722}.

\bibitem{Ozawa}
N.~Ozawa, \emph{About the QWEP conjecture},
Internat. J. Math. \textbf{15} (2004), 501--530.
\doi{10.1142/S0129167X04002417}.

\bibitem{Pestov}
V.~G. Pestov, \emph{Hyperlinear and sofic groups: a brief guide},
Bull. Symbolic Logic \textbf{14} (2008), 449--480.
\doi{10.2178/bsl/1231081461}.

\bibitem{Steinberg}
B.~Steinberg, \emph{Simplicity, primitivity and semiprimitivity of
\'etale groupoid algebras with applications to inverse semigroup algebras},
\href{https://arxiv.org/abs/1408.6014}{arXiv:1408.6014} (2014).

\bibitem{Stepanov}
A.~V. Stepanov, \emph{On the normal structure of the general linear
group over a ring}, Zap. Nauchn. Sem. POMI \textbf{236} (1997), 166--182;
English transl., J. Math. Sci. \textbf{95} (1999), 2146--2155.
\doi{10.1007/BF02169976}.

\bibitem{Thom}
A.~Thom, \emph{Examples of hyperlinear groups without factorization
property}, Groups Geom. Dyn. \textbf{4} (2010), 195--208.
\doi{10.4171/GGD/80}.

\end{thebibliography}
\end{document}
